\documentclass[11pt]{article}
\usepackage[a4paper,margin=26mm]{geometry}
\usepackage{amsmath,amssymb,amsthm,booktabs,longtable,array}
\usepackage[unicode,hidelinks]{hyperref}
\usepackage{fontspec}
\setlength{\emergencystretch}{2em}
\setmainfont{lmroman10-regular.otf}[BoldFont=lmroman10-bold.otf,ItalicFont=lmroman10-italic.otf,BoldItalicFont=lmroman10-bolditalic.otf]
\newtheorem{theorem}{Theorem}
\newtheorem{proposition}[theorem]{Proposition}
\newtheorem{lemma}[theorem]{Lemma}
\newtheorem{corollary}[theorem]{Corollary}
\theoremstyle{remark}\newtheorem{remark}[theorem]{Remark}
\newcommand{\Aut}{\operatorname{Aut}}
\newcommand{\ord}{\operatorname{ord}}
\newcommand{\PSL}{\operatorname{PSL}}
\newcommand{\PSU}{\operatorname{PSU}}
\newcommand{\GL}{\operatorname{GL}}
\newcommand{\SL}{\operatorname{SL}}
\newcommand{\GU}{\operatorname{GU}}
\newcommand{\SU}{\operatorname{SU}}
\title{Chiral polytopes by prefix coordinates\\\large A complete finite enumeration contract}
\author{human\_compression research project}
\date{11 October 2026 --- author manuscript}
\begin{document}\maketitle
\begin{abstract}
We recover the finite chiral-polytope enumeration of Buekenhout, Leemans and Tranchida through a shared theory of marked coordinates and exact orbit searches. The subgroup dictionary is literal: interval-rotation subgroups become difference subgroups without changing any intersection. We prove termination, completeness, the reflection test, and both oriented and unrooted equivalence conventions. Exact small-group enumerations validate the whole search, while uniform counts reconstruct the mathematical columns of the source timing table. The geometric reconstruction theorem of Schulte and Weiss and the almost-simple involution bound of Liebeck and Shalev remain attributed inputs. The source's runtime measurements are retained as empirical observations. The distinction between directly regular and nonorientable regular output is made explicit.
\end{abstract}

\section{Source, scope and geometric input}

Our principal source is the complete fourteen-page revised author manuscript
\emph{A new algorithm to classify chiral polytopes with a given automorphism group}, by Francis Buekenhout, Dimitri Leemans and Philippe Tranchida, deposited as \texttt{algo-rev.pdf} at the Universit\'e libre de Bruxelles. We denote it by BLT. Its journal citation is \emph{Beitr\"age zur Algebra und Geometrie} \textbf{62} (2021), 21--36. The source ledger identifies the precise PDF bytes and separately records the earlier arXiv:1912.02287v1. The journal's sixteen-page typesetting has not been obtained, so a statement about the revised author text is not presented as a verified statement about every journal page.

The principal mathematical contract comprises BLT Theorem 4.3, Lemmas 4.4--4.5, and the complete search described in Section 4.2. We also recover the mathematical candidate-space comparison in Section 4.3, retaining its external involution bound. The only numerical table in BLT contains timings, group orders and involution proportions. It is not a table of numbers of chiral polytopes. Our exact small-group census is a validation of the recovered algorithm, not an invented reconstruction of an absent source census.

Throughout, isomorphisms preserve face ranks. Duality is a different equivalence and is not factored out. A connected Schl\"afli diagram means that every adjacent entry of its string symbol exceeds two. This is the nondegenerate scope stipulated in BLT Section 4. An \emph{oriented} rotary polytope is supplied with a choice of one of the two parity classes of flags. Removing this choice gives the underlying unrooted polytope. In a chiral polytope these two choices are inequivalent under automorphisms of the polytope, but they have the same underlying partially ordered set.

We use the following classical reconstruction theorem, with its original attribution. It is Schulte--Weiss, \emph{Chiral polytopes} (1991), Theorem 1, in the form restated in BLT Theorem 4.1 and Hartley--Hubard--Leemans (2012), printed pp.374--375. Let
\[
G=\langle s_1,\ldots,s_{n-1}\rangle,\qquad n\geq3,
\]
where $\ord(s_i)=p_i\geq3$ and every consecutive product
$s_i\cdots s_j$ with $i<j$ is a nonidentity involution. For
$J\subseteq\{0,\ldots,n-1\}$ define
\[
T_J=\langle s_{i+1}\cdots s_j:i<j,\ i,j\in J\rangle.
\tag{1}
\]
Suppose $T_J\cap T_K=T_{J\cap K}$ for all $J,K$. There is then a chiral or directly regular polytope of type
$\{p_1,\ldots,p_{n-1}\}$ with rotation group $G$. It is directly regular precisely when there is an automorphism $f$ satisfying
\[
f(s_1)=s_1^{-1},\quad f(s_2)=s_1^2s_2,\quad f(s_i)=s_i\ (i\geq3).
\tag{2}
\]
Its chain stabilizers in the rotation group are the groups in (1). Conversely, a chiral polytope with a chosen flag yields this marked data. These geometric existence and stabilizer assertions are retained, not reproved under the name of coordinate transport.

The exact-order formulation above uses a slightly stronger explicit input than the usual presentation by squared interval words. It holds for actual polytopes: the corresponding two distinct flag adjacencies cannot have product equal to the identity. It is also what the recovered enumeration tests. There is no silent assumption that arbitrary groups satisfying only the string relations are polytopal.

\section{The complete subgroup dictionary}

Put $a_0=1$ and $a_i=s_1\cdots s_i$. For $J\subseteq\{0,\ldots,n-1\}$ put
\[
D_J=\langle a_i^{-1}a_j:i,j\in J\rangle.
\tag{3}
\]
The shared marked-coordinate theorem MC.1 gives, simultaneously for every $J$,
\[
s_i=a_{i-1}^{-1}a_i,\qquad
s_{i+1}\cdots s_j=a_i^{-1}a_j,\qquad T_J=D_J.
\tag{4}
\]
Thus generation and the entire intersection diagram are preserved, rather than merely the top generated group. The empty and singleton subgroups in (3) are trivial; consequently the intersection property is automatic if either indexing set has at most one member. Adjoining the formal improper indices $-1,n$ with their associated trivial words changes no subgroup.

\begin{theorem}[BLT prefix characterization, in ordered coordinates]
Let $G$ be finite. Its connected oriented chiral polytopes are represented by ordered normalized tuples
$\boldsymbol a=(1,a_1,\ldots,a_{n-1})$ satisfying all of the following conditions:
\begin{enumerate}
\item $n\geq3$ and $\langle a_1,\ldots,a_{n-1}\rangle=G$;
\item $\ord(a_{i-1}^{-1}a_i)>2$ for $1\leq i<n$;
\item $\ord(a_i^{-1}a_j)=2$ whenever $0\leq i<j<n$ and $j-i\geq2$;
\item $D_J\cap D_K=D_{J\cap K}$ for all $J,K$;
\item no automorphism of $G$ maps every $a_i$ to $a_i^{-1}$.
\end{enumerate}
Two tuples represent the same oriented isomorphism class exactly when an automorphism of $G$ sends the entries of one to the corresponding entries of the other.
\end{theorem}

\begin{proof}
By (4), conditions 1--4 are exactly the generating, order and intersection conditions in the retained reconstruction theorem. In particular $a_i$ is an involution for $i\geq2$, while $a_1$ has order greater than two. MC.5 identifies the automorphism in (2) with simultaneous inversion of the $a_i$. Such an automorphism has square equal to the identity because its square fixes a generating set, and it is nonidentity because it inverts $a_1$. Thus condition 5 is precisely chirality.

A generating-marking isomorphism transports every subgroup (3), and hence the corresponding geometric construction and its chosen flag class. Conversely, an orientation-preserving polytope isomorphism may be followed by an automorphism in the transitive rotation group so that it sends a chosen base flag to the other chosen base flag. Its conjugation action then sends each distinguished rotation, and hence each prefix, to the corresponding one. This proves the exact oriented equivalence. An explicit flag-graph proof of the same statement is given in the next section.
\end{proof}

For completeness, the face-stabilizer formulas in BLT Proposition 3.4 follow directly from (1). Omitting index zero leaves
$\langle s_2,\ldots,s_{n-1}\rangle$, and omitting index $n-1$ leaves
$\langle s_1,\ldots,s_{n-2}\rangle$. Omitting an internal index $i$ leaves
\[
\left\langle\{s_j:j\ne i,i+1\}\cup\{s_is_{i+1}\}\right\rangle.
\tag{5}
\]
Every interval word avoiding endpoint $i$ either avoids the crossing or crosses it using the paired factor $s_is_{i+1}$, proving both inclusions in (5). Restriction to a section retains the corresponding contiguous rotation and subgroup data. The retained reconstruction theorem therefore gives the directly regular-or-chiral section statement of BLT Proposition 3.3. The fundamental connection between flag stabilizers and these subgroups remains the classical input.

\section{Flag graphs, reflection and the missing mirror quotient}

The flag graph of the retained construction has $2|G|$ vertices. A useful convention labels them $G\times\{+,-\}$ and sets
\[
 r_i(g,+)=(ga_i,-),\qquad r_i(g,-)=(ga_i^{-1},+)
 \quad(0\leq i<n).
\tag{6}
\]
This convention interchanges the two parity labels relative to the convention in which $a_i$ is the right-action word $\rho_0\rho_i$. To check the identification, represent the two cosets of the universal even subgroup by $w$ and $w\rho_0$. Right multiplication by $\rho_i$ sends these to
$w a_i^{-1}\rho_0$ and $w a_i$, respectively. Exchanging the displayed parity labels gives (6). Thus (6) is the actual flag graph furnished by reconstruction, not an unrelated graph with the same order.

Each $r_i$ is an involution without fixed points. Distant $r_i,r_j$ commute because $(a_i^{-1}a_j)^2=1$. The graph is connected since the tuple generates $G$. Left multiplication by $G$ acts on it, freely and transitively on each parity class.

\begin{proposition}[Exact flag isomorphisms]
Let $\boldsymbol a,\boldsymbol b$ be valid tuples of the same rank in finite groups of the same order. Every color-preserving isomorphism between their flag graphs, after a left translation in the target, is one of the following:
\begin{align*}
 (g,\epsilon)&\longmapsto(f(g),\epsilon),
 & f(a_i)&=b_i;\\
 (g,\epsilon)&\longmapsto(f(g),-\epsilon),
 & f(a_i)&=b_i^{-1},
\end{align*}
where $f$ is a group isomorphism. Consequently, unrooted chiral polytopes are classified by the automorphism orbits of tuples followed by the extra identification
\[
 (a_1,\ldots,a_{n-1})\sim(a_1^{-1},\ldots,a_{n-1}^{-1}).
 \tag{7}
\]
Every unrooted chiral class consists of exactly two oriented classes.
\end{proposition}

\begin{proof}
A connected bipartite graph has just two choices of its ordered bipartition. Hence an isomorphism either preserves or exchanges the displayed parity classes. Translate its image of $(1,+)$ to the appropriate vertex with first coordinate 1. Since color zero exchanges $(g,+)$ and $(g,-)$, the same map $f$ occurs in both coordinates. For a parity-preserving map, preservation of color $i$ gives
$f(ga_i)=f(g)b_i$. For a parity-exchanging map it gives
$f(ga_i)=f(g)b_i^{-1}$. MC.3, or its word argument, shows that $f$ is the required group isomorphism. Conversely these equations verify every colored edge directly.

In particular, switching the parity labels is an isomorphism from the graph of $\boldsymbol a$ to the graph of $\boldsymbol a^{-1}$. The corresponding polytopes have the same unrooted isomorphism class. The two oriented classes coincide exactly when an automorphism of $G$ inverts all entries. This is exactly the regularity condition excluded in Theorem 1. Therefore the mirror action is free on the chiral oriented classes, and each unrooted class has two of them.
\end{proof}

A color-preserving isomorphism of polytope flag graphs gives an isomorphism of the polytopes: an $i$-face is the connected component obtained by deleting color $i$, and incidence is witnessed by a common flag. Thus the graph calculation proves the claimed unrooted equivalence, not only a graph equivalence. It agrees with the explicit discussion of enantiomorphic base flags in Hartley--Hubard--Leemans, printed p.374.

The distinction has an immediate computational consequence. The final deduplication in the printed BLT code tests only a homomorphism sending a tuple to another tuple. It naturally identifies oriented classes. If ``up to isomorphism'' means unrooted poset isomorphism as defined in the source's Section 2, the additional step (7) is necessary. We provide both outputs. This is a convention correction with a reproducible finite witness, not a claim that the source's choice of prefixes loses chiral polytopes.

\section{An exhaustive search in every rank}

Here is a fully specified algorithm, independent of a computer algebra system's choice of generating sequence or conjugacy-class ordering. The input is a finite group with exact multiplication, inverse and equality operations.
\begin{enumerate}
\item List all nonidentity involutions $\mathcal I$. Choose one representative $a_1$ of every conjugacy class of elements of order greater than two.
\item At a prefix $(1,a_1,\ldots,a_k)$ compute its subgroup and all its difference subgroups (3), using the star formula MC.2.
\item If the prefix generates $G$, record it when $k\geq2$ and do not extend it. Otherwise try every $x\in\mathcal I$ outside its generated subgroup. Require
\[
 \ord(a_k^{-1}x)>2,\qquad
 \ord(a_i^{-1}x)=2\quad(0\leq i<k),
 \tag{8}
\]
and test the full difference-subgroup intersection property on the enlarged prefix. Recurse exactly for the candidates that pass.
\item For each generating tuple, use MC.3 to test whether its simultaneous inverse tuple is isomorphic to it. Retain the tuples for which this test fails.
\item Remove duplicates using the canonical colored Cayley array of MC.3. This gives the oriented list. Pair the resulting list by (7) for the unrooted list. For either list report
$\{\ord(a_{i-1}^{-1}a_i):1\leq i<n\}$ in order.
\end{enumerate}

\begin{theorem}[Complete finite enumeration]
The algorithm terminates for every finite input group. Its two outputs contain exactly one representative of every connected oriented, respectively unrooted, chiral polytope with full automorphism group isomorphic to the input group. No rank cutoff or list of possible Schl\"afli types is required.
\end{theorem}

\begin{proof}
Every initial subtuple of a valid tuple satisfies the same order and difference-subgroup intersection tests in its generated subgroup. These properties are invariant under conjugation, so normalizing the first entry loses no tuple. The intersection identity applied to $\{0,\ldots,k\}$ and $\{0,k+1\}$ shows that a valid next entry cannot belong to the old generated subgroup. It follows both that our exclusion of such entries is safe and that generation cannot occur before the end of a valid tuple.

MC.4 now proves termination and exhaustiveness of the prefix search. More explicitly, each extension at least doubles the order of the prefix subgroup and the first subgroup has order at least three. A rank-$n$ recorded tuple therefore satisfies $3\cdot2^{n-2}\leq |G|$. The search tests all permitted new entries at every stage, so every normalized valid tuple occurs along one of its branches. Theorem 1 proves that the reflection filter retains exactly the chiral cases and that their full automorphism group is $G$. MC.3 gives precisely the oriented equivalence; Proposition 2 gives precisely the extra unrooted equivalence. This also proves the correctness of the reported types.
\end{proof}

The normalizer and centralizer pruning used by BLT follows directly from (8). A nonadjacent involutory entry $x$ relative to $a_1$ satisfies
$(a_1^{-1}x)^2=1$, equivalently $xa_1x=a_1^{-1}$, so it lies in the normalizer of $\langle a_1\rangle$. With an earlier involution $a_i$, the condition $(a_ix)^2=1$ is equivalent to commutation; the order is exactly two because the entries are distinct. For the immediately previous involution, order greater than two is equivalent to noncommutation. Restrictions inherited from all earlier nonadjacent positions must be retained when the next entry is tried; the restriction from the preceding adjacent position must not be incorrectly retained after it becomes nonadjacent. Direct order tests avoid this bookkeeping issue. Stabilizer-orbit pruning is optional and is justified by U.ExtensionOrbits as incorporated in the companion paper, Section 2.

Our implementation uses all candidates after the first-entry conjugacy reduction. It does not rely on the order in which software lists conjugacy classes. It tests the first adjacent order explicitly, and it uses the supplied ordered tuple rather than asking whether an implicit computer algebra generating sequence happens to have the intended order. These choices produce a transparent reference implementation; they are not a claim to improve the original optimized runtime.

\section{Candidate counts and the source timing table}

Write $N=|G|$, $v=|\mathcal I|$, and $u=N-1-v$, the number of elements of order greater than two. Before subgroup, order or conjugacy pruning, the original rank-$n$ rotation search has at most $u^{n-1}$ tuples, while the prefix search has at most
\[
 u\,v^{n-2}.
 \tag{9}
\]
Indeed only its first nonidentity coordinate is selected from the $u$ noninvolutions, and all remaining coordinates are involutions. If $u>0$, the ratio of these Cartesian candidate counts is $(v/u)^{n-2}$. Distinctness, conjugacy representatives and prefix rejection can only reduce the counts. Equation (9) compares candidate sets; it is not a worst-case ratio between complete running times, since individual subgroup and isomorphism tests also have costs.

BLT Proposition 4.6 retains the Liebeck--Shalev involution estimate: for almost simple groups outside the family with socle $\PSL_2(q)$, there is an absolute constant $c$ with $v<cN^{5/8}$. We retain precisely this attributed input as quoted in BLT; its classification-theoretic proof is not included in either of our two papers. It implies $v/N=O(N^{-3/8})$, and hence $v/u=O(N^{-3/8})$ along such a family. Thus (9) explains the decreasing candidate proportion at each fixed rank. No uniform small numerical value is implied for every small group by an unspecified constant.

The timings in the source were measured on the stated machine with Magma. They are empirical records and cannot follow as corollaries of an abstract proof. We do not recreate their hardware or claim to reproduce those timings. Their group-order and involution-proportion columns, however, have exact mathematical content. The following uniform formulas reconstruct them.

\begin{proposition}[Exact involution counts for all table families]
Let $q$ be a prime power. Then
\begin{align*}
 |\PSL_3(q)|&=\frac{q^3(q^3-1)(q^2-1)}{(3,q-1)},&
 i_2(\PSL_3(q))&=\begin{cases}q^2(q^2+q+1)&q\text{ odd},\\(q^3-1)(q+1)&q\text{ even};\end{cases}\\
 |\PSU_3(q)|&=\frac{q^3(q^3+1)(q^2-1)}{(3,q+1)},&
 i_2(\PSU_3(q))&=\begin{cases}q^2(q^2-q+1)&q\text{ odd},\\(q^3+1)(q-1)&q\text{ even}.
 \end{cases}
\end{align*}
Also
\[
 i_2(S_n)=\sum_{k=1}^{\lfloor n/2\rfloor}\frac{n!}{2^k k!(n-2k)!},\qquad
 i_2(A_n)=\sum_{\substack{1\leq k\leq n/2\\k\text{ even}}}\frac{n!}{2^k k!(n-2k)!}.
 \tag{10}
\]
\end{proposition}

\begin{proof}
Counting ordered bases gives
$|\GL_3(q)|=(q^3-1)(q^3-q)(q^3-q^2)$. Divide by the $q-1$ determinant values and then by the center of order $(3,q-1)$ to obtain the first group order. In odd characteristic an involution in $\SL_3(q)$ has a one-dimensional $+1$ eigenspace and a complementary two-dimensional $-1$ eigenspace. There are $q^2+q+1$ choices of the line and $q^2$ complements of it. Conversely every such ordered decomposition determines one involution. In characteristic two, an involution is $1+u\varphi$, where $u\ne0$, $\varphi\ne0$, and $\varphi(u)=0$. Its nilpotent part has rank one because its square is zero in dimension three. For each image line, fix one nonzero $u$; there are $q^2-1$ suitable nonzero functionals. Thus the count is $(q^2+q+1)(q^2-1)=(q^3-1)(q+1)$.

For the unitary calculation use a nondegenerate Hermitian form on $\mathbb F_{q^2}^d$, and let $z_d$ be the number of norm-zero vectors. Each nonzero value of the field norm has $q+1$ preimages, while zero has one. Hence
$z_{d+1}=(q+1)q^{2d}-qz_d$, starting with $z_1=1$. This gives
$z_2=q^3+q^2-q$ and $z_3=q^5-q^3+q^2$. There are consequently $q^3+1$ isotropic lines and $q^2(q^2-q+1)$ nonsingular lines in dimension three. The counts of unit vectors in dimensions 1, 2 and 3 are, respectively,
$q+1$, $q(q^2-1)$ and $q^2(q^3+1)$. Counting orthonormal bases gives
$|\GU_3(q)|=q^3(q^3+1)(q^2-1)(q+1)$; the determinant map has $q+1$ values, giving the stated order of $\SU_3(q)$ before quotienting its center.

In odd characteristic, an involution in $\SU_3(q)$ is $+1$ on a nonsingular line and $-1$ on its orthogonal plane. Thus the number is exactly the number of nonsingular lines. In characteristic two its nilpotent part again has rank one. If its image is $\langle u\rangle$, preservation of the form implies that its kernel is $u^\perp$. The inclusion of image in kernel makes $u$ isotropic. It therefore has the form
$N(x)=c\,h(x,u)u$, with $c+c^q=0$. Conversely this formula gives a determinant-one unitary involution for every nonzero such $c$. For a fixed nonzero representative $u$ of an isotropic line there are $q-1$ choices, giving $(q^3+1)(q-1)$.

In both special groups the center has odd order dividing three. Every projective involution has a unique involutory lift: if $x^2$ is central, multiplication by the unique central square root of $x^{-2}$ produces the lift. The projective counts therefore equal the counts just obtained. Finally a permutation involution consists of $k$ disjoint transpositions; selecting and pairing their $2k$ moved points gives the summand in (10). It is even precisely when $k$ is even.
\end{proof}

The accompanying exact table records all 32 rows of the revised author source. Twenty of the counts are checked independently by GAP conjugacy classes: both degree-three classical families for $q=2,3,4,5$, and both symmetric and alternating families for $5\leq n\leq10$. The all-parameter conclusion is the written proof, not an extrapolation from these twenty checks. For example, the revised source's unitary ratios $0.053$ for $\PSU_3(2)$ and $0.0097$ for $\PSU_3(4)$ differ from the exact values $9/72=1/8$ and $195/62400=1/320$. Its $\PSU_3(5)$ ratio is $525/126000=1/240$. The revised source already corrects the earlier arXiv entry for $\PSL_3(17)$; that earlier error is not attributed to the revised text.

\section{Exact models and the regular-output boundary}

The reference enumeration includes the trivial group and elementary abelian groups of orders two and four as empty-output boundary cases. It constructs finite multiplication tables from explicit permutations and searches every rank allowed by the subgroup-growth proof. Its counts are as follows. The regular column counts connected directly regular polytopes with the indicated group as rotation group, with duals kept distinct.
\begin{center}
\begin{tabular}{lrrr}\toprule
Input group&Oriented chiral&Unrooted chiral&Directly regular\\\midrule
$C_5$&0&0&0\\
$D_{10}$&0&0&0\\
$A_4$&0&0&1\\
$S_4$&0&0&2\\
$A_5$&0&0&4\\
$\operatorname{AGL}_1(5)$&2&1&0\\
$\operatorname{AGL}_1(7)$&4&2&0\\
$\PSL_2(7)$&0&0&5\\\bottomrule
\end{tabular}
\end{center}
For $A_5$ the four types are $\{3,5\},\{5,3\},\{5,5\}$ and the rank-four type $\{3,3,3\}$. For $\operatorname{AGL}_1(5)$ the chiral type is $\{4,4\}$; for $\operatorname{AGL}_1(7)$ they are $\{3,6\}$ and $\{6,3\}$. An independent rank-three search tries every ordered pair of original rotations, verifies generation and all three cyclic subgroup intersections, and agrees with every rank-three prefix output. The saved witnesses also undergo a complete color-preserving graph-automorphism search and the component-intersection tests at both flag orbits. The mirror isomorphism in (7) is checked on every colored edge.

The group of order 20 can be written
$\mathbb F_5\rtimes C_4$, with multiplication
$(u,i)(v,j)=(u+2^iv,i+j)$.
Take $a_1=(0,1)$ and $a_2=(1,2)$. Their successive rotation orders are 4 and 4, and the three relevant cyclic subgroups have pairwise trivial intersections. They generate the whole group. Its characteristic derived subgroup is the translations. An automorphism conjugates the action of the quotient on this subgroup through an automorphism of $C_5$; that automorphism group is abelian, so the quotient action $u\mapsto2u$ cannot be changed to $u\mapsto2^{-1}u$. Thus no automorphism inverts $a_1$. This gives the explicit chiral mirror-pair example without relying solely on a census count.

There are two further boundary controls. First, the tetrahedral $A_4$ marking is inverted by an outer automorphism, but not by any inner automorphism. Replacing $\Aut(G)$ by $G$ in the reflection test would falsely declare it chiral. This is also why the automorphism-language repair already present in the revised BLT Lemma 4.4 must be preserved.

Second, the final paragraph of BLT says that removing the chirality test also computes all regular polytopes when $G$ is simple. The retained reconstruction theorem instead returns directly regular polytopes with rotation group $G$; these have $2|G|$ flags and full group of order $2|G|$. A regular polytope whose full group is the nonabelian simple group $G$ has $|G|$ flags. Its rotation subgroup is $G$ because $G$ has no subgroup of index two, so it is nonorientable and is not the same reconstructed object.

An explicit certificate uses $A_5$ on $\{0,1,2,3,4\}$ with
\[
r_0=(1\,3)(2\,4),\qquad
r_1=(0\,3)(2\,4),\qquad
r_2=(1\,4)(2\,3).
\tag{11}
\]
These generate $A_5$, the end generators commute, their consecutive product orders are 3 and 5, and all 64 standard subgroup-intersection tests pass. Thus they define a regular 60-flag polytope. Their prefix rotation tuple
$(1,r_0r_1,r_0r_2)$ generates $A_5$ and passes the difference-subgroup tests. Applying the rotary reconstruction gives a directly regular 120-flag polytope. The two flag counts exclude isomorphism. Consequently the last paragraph needs an additional nonorientable quotient or an independent ordinary C-group search; it does not follow by simply deleting the filter. The chiral enumeration theorem is unaffected.

\section{Source clauses, retained costs and verification limits}

The source ledger assigns each principal and supporting clause to a specific proof span in this TeX source or in \emph{Marked coordinates and exact orbit enumeration}. The latter paper is a reusable acquisition; this application pays separately for the geometric interface, the actual admissibility and equivalence arguments, and the exact table calculations. The source's use of the Schulte--Weiss theorem and the Liebeck--Shalev bound is retained with attribution. The prefix idea, the predominance of involutory candidates and the reflection test are the source authors' contributions and inputs, not credited to the model as inventions.

The revised text still has transcription inconsistencies. Its equation (6) includes $i=j$ in the squared interval relation, which would force each adjacent rotation to have order at most two; it must read $i<j$ in this use. Equation (10) gives distance two in both cases; the noninvolutory case must have distance one. Theorem 1 states the consistent conditions explicitly. The index formulas in its discussion of $A_I$ and the piecewise expansion inside Theorem 4.3 are replaced by the exact dictionary (4). We distinguish these corrections from the already corrected automorphism language and $\PSL_3(17)$ entry in the revised edition. The omitted explicit first adjacent-order check in the printed recursion is not used to assert that a chiral output is necessarily wrong; our implementation simply tests every condition in its declared contract.

The seven Lean declarations in \texttt{PrefixCoordinates.lean} verify the all-length coordinate cancellation, its two inverses, the equality of actual subgroup closures, intersection transport, the reflected second generator, and the involutivity of a simultaneous-inversion automorphism on a generating set. They use only standard reported axioms. They do not formalize the Schulte--Weiss polytope construction, the whole search, the Lie-theoretic input bound, or the classification table calculations. The finite scripts and negative controls supplement the written proof; they do not replace its arbitrary-group and arbitrary-rank quantifiers.

No new conjecture is stated in this source's algorithm sections, and no open-question solution is claimed by this recovery. The finite corrections and equivalence clarification are kept in the source contract rather than assigned an artificial infinite compression ratio. A later matched measurement must charge the full relevant umbrella and the full retained-input boundary, together with the application. This paper makes no numerical compression or performance claim.

\begin{thebibliography}{9}
\bibitem{BLT} F.~Buekenhout, D.~Leemans and P.~Tranchida, \emph{A new algorithm to classify chiral polytopes with a given automorphism group}, Beitr. Algebra Geom. \textbf{62} (2021), 21--36. Exact principal edition: ULB revised author manuscript, \url{https://dipot.ulb.ac.be/dspace/bitstream/2013/315729/3/algo-rev.pdf}; comparison edition: arXiv:1912.02287v1. DOI: 10.1007/s13366-020-00536-1.
\bibitem{HHL} M.~I.~Hartley, I.~Hubard and D.~Leemans, \emph{Two atlases of abstract chiral polytopes for small groups}, Ars Math. Contemp. \textbf{5} (2012), 371--382; especially pp.374--375 for the retained reconstruction and enantiomorphic base flags.
\bibitem{SW} E.~Schulte and A.~Ivi\'c Weiss, \emph{Chiral polytopes}, DIMACS Ser. Discrete Math. Theoret. Comput. Sci. \textbf{4} (1991), 493--516, Theorem 1. The exact retained statement is read in BLT and HHL; the full original chapter was not obtained in this audit.
\bibitem{LS} M.~W.~Liebeck and A.~Shalev, \emph{Classical groups, probabilistic methods, and the $(2,3)$-generation problem}, Ann. of Math. \textbf{144} (1996), 77--125, Corollary 4.5 as retained in BLT Proposition 4.6. Its proof is not claimed recovered here.
\bibitem{FLW} M.~E.~Fernandes, D.~Leemans and A.~Ivi\'c Weiss, \emph{Highly symmetric hypertopes}, Aequationes Math. \textbf{90} (2016), 1045--1067. This is the source of the $C^+$-group viewpoint used by BLT; our string application uses the classical polytope reconstruction directly.
\bibitem{MC} human\_compression research project, \emph{Marked coordinates and exact orbit enumeration}, companion TeX manuscript, phase 282. Its orbit interface reuses U.ExtensionOrbits from the shared proof bank and accepted enumeration applications.
\end{thebibliography}
\end{document}
